\title{Zinc doping of the Cuprates}
\author{Henry Sheehy}
\date{ \today }
  
\maketitle
  
\tableofcontents

\section{Introduction}
J.C. Seamus Davis discussed with me the problem of zinc doping in the cuprates. 
It unknown why the superconductivity is killed by small doping.
Seamus suggested I try modelling the situation by switching off the Hubbard U on the zinc site.
The cuprate can be captured by a stacked SSH model, similar to that studied by Rosenberg and Manousakis [CITE: Rosenberg, Aryal PRB 100 104522 (2019)].

H. Alloul {\it et al} review Defects in correlated metals and superconductors [CITE:arXiv:0711.0877v1 (2018)].

\section{Topology}
Rosenberg and Manousakis [CITE: Rosenberg, Aryal PRB 100 104522 (2019)] showed topological modes when the interdimer hopping is different from the intradimer hopping, with equal chemical potential.

{\color{red} Have they checked with a rigid shift?}

{\color{red} Any nonzero CDW-like density in the Fock channel would give rise to superconductivity.}


\section{Review of defects in correlated metals}
H. Alloul {\it et al} [CITE:arXiv:0711.0877v1 (2018)] review  theories of impurities in $d$-wave superconductor theories of impurities 


\section{Reproducing results}
The normal state Hamiltonian is given by
\begin{multline}
    \hat{\mathcal{H}}_{ 0 } =
        \sum_{0\leq x,y\leq N}\sum_\sigma\sum_{m\in\{A,B\}]}\left[
            -\mu_m \hat{c}^{ \sigma,m \dagger}_{ x,y }\hat{c}^{ \sigma,m }_{ x,y }
        \right]\\
        -t \sum_{0\leq x,y\leq N}\sum_{\sigma}\left[
            \hat{c}^{ \sigma,A \dagger}_{ x,y }\hat{c}^{ \sigma,B }_{ x,y }
            +
            \hat{c}^{ \sigma,A \dagger}_{ x,y }\hat{c}^{ \sigma,B }_{ x,y+1 }
            +
            \hat{c}^{ \sigma,B \dagger}_{ x,y }\hat{c}^{ \sigma,A }_{ x+1,y+1 }
            +
            \hat{c}^{ \sigma,B \dagger}_{ x,y }\hat{c}^{ \sigma,A }_{ x+1,y }
            +\text{h.c.}
        \right]
\end{multline}
where the fermion operator $\hat{c}^{\sigma,m\dagger}_{ x,y }$ creates an operator on atom $m\in\{A,B\}$ within the unit cell located at $\mathbf{R}_{x,y}= x \mathbf{b}_0\hat{x}+y \mathbf{b}_1 \hat{y}$,
where $\mathbf{b}_0, \mathbf{b}_1$ are the lattice basis vectors,
$\hat{n}^{(\alpha)\dagger}_{ x,y }$ is the number operator, $\mu$ is the chemical potential and $s$ is a rigid shift of the atomic potentials.
The intra- and inter-unit-cell hoppings strengths in the $\hat{x}$ direction are given by $v$ and $w$, which form the terms of the standard SSH model.
These 1D SSH chains are coupled by a diagonal hopping $t_d$, which form a two-dimensional stacking.
The lattice is depicted with $N=3$ unit cells in Figure \ref{fig:2D-SSH_unit_cell}. 
Following [CITE: Rosenberg and Manousakis arXiv:2203.12004v1], without loss of generality, we set $v=0.6, w=1.2$, the topological phase of the SSH model, and $t_d=0.9$.

\createfigure{2D-Weyl-SSH/unit_cell}{2D-SSH_unit_cell}

The solution to the Hamiltonian is given by the Green's function, which we remind the reader, is given by
\begin{equation}
    \hat{\mathcal{G}}(\omega-i\epsilon,x,y)
    \equiv
    \sum_n\left[
        \frac{\psi_n(x,y)\psi_n(x,y)^\dagger}{\omega-i\epsilon-E_n}
        \right],
\end{equation}
where $\psi_n(x,y)$ is the eigenvector associated with the $n^\text{th}$ eigenenergy, resolved at the point on the lattice $(x,y)$. 
The Green's function matrix has a component for each atom.
The Green's function is experimental observables through the {\it local density of states},
\begin{equation}
    \text{LDOS}(\omega;x,y)
    \equiv
    -\frac{1}{\pi}\Im\hat{\mathcal{G}}(\omega-i\epsilon,x,y)
,
\end{equation}
which is a measure of the density of states at a given point, at a specified applied bias energy $\omega$ and resolution $\epsilon$.
The local density of states for each atomic site with the unit cells are given in Figure \ref{fig:2D-SSH_ldos_each_atom}, showing the existence of zero energy edge states.

\createfigure{2D-Weyl-SSH/ldos_each_atom}{2D-SSH_ldos_each_atom}

Resolving the local density states at the level of the unit cell (summing over the atoms) and integrating over the $y$-axis, we study the {\it spectrum} of the system. 
The spectrum shows edge modes present at zero energy at the edges, which are not present in the bulk. These are distinct signatures of Majorana modes. 
In addition to localised density at the edge, there exists a `tail' of density which the authors investigate further.

\createfigure{2D-Weyl-SSH/ldos}{2D-SSH_ldos}

\createfigure{2D-Weyl-SSH/spectrum}{2D-SSH_spectrum}
Majorana modes exist pinned to {\it Dirac points}, which are the cross point of two bands. 
Studying the band structure in Figure \ref{fig:2D-SSH_k_space}, we observe a zero-mode between the Dirac points, which is known as a {\it Bogoliubov-Fermi arc}. 

\createfigure{2D-Weyl-SSH/k_space}{2D-SSH_k_space}

In order to determine whether the modes are topological modes, it must be shown that they their wavefunctions are distributed nonlocally across the system.
We plot the wavefunction via the related {\it spectral density} of the Green's function at the point $k_y$ at which the Bogoliubov-Fermi arc and Majorana modes are are located, Figure \ref{fig:majorana_fermi_arc}.
As a function of $x$-position, we observe that indeed, the modes are spread non-locally and that the arc tails.

\createfigure{2D-Weyl-SSH/majorana_fermi_arc}{majorana_fermi_arc}

We note that other hopping configurations are possible, leading either to topological modes along one pair of parallel edges, or along all edges; see 
[Two-Dimensional Extended Su-Schrieffer-Heeger Model; Bo Hung; Masters Thesis; National Taiwan Normal University] 
.

\section{Interactions and mean-field solution}
We add an interaction term to model multiorbital superconductivity within the unit-cell, $U_v$, and an interorbital intercell Coulomb repulsion term, $U_w$. The interacting contribution to the Hamiltonian is 
\begin{equation}
    \hat{\mathcal{H}}_{ I }=
    -U_v \hat{c}^{ (A) \dagger}_{ x,y }\hat{c}^{ (B) \dagger}_{ x,y } \hat{c}^{ (A) }_{ x,y }\hat{c}^{ (B) }_{ x,y }
    +U_w \hat{c}^{ (A) \dagger }_{ x,y }\hat{c}^{ (B) \dagger}_{ x-1,y } \hat{c}^{ (B) }_{ x-1,y }\hat{c}^{ (A) }_{ x,y }
.
\end{equation}
The mean-field Hamiltonian is
\begin{multline}
    \hat{\mathcal{H}}^{'}_{ I }=
    \sum_{x,y}\Bigg[
    \sum_m\left[\phi_{\,m}\hat{n}^{ (m) }_{ x,y }\right]
    +\Phi_v \hat{c}^{ (A) }_{ x,y }\hat{c}^{ (B) }_{ x,y }
    +\Phi_w \hat{c}^{ (A) }_{ x,y }\hat{c}^{ (B) }_{ x-1,y }
    \\
    +\Delta_v \hat{c}^{ (A) \dagger}_{ x,y }\hat{c}^{ (B) \dagger }_{ x,y }
    +\Delta_w \hat{c}^{ (A) \dagger}_{ x,y }\hat{c}^{ (B) \dagger }_{ x-1,y }
    +\Delta_v^* \hat{c}^{ (A) }_{ x,y }\hat{c}^{ (B) }_{ x,y }
    +\Delta_w^* \hat{c}^{ (A) }_{ x,y }\hat{c}^{ (B) }_{ x-1,y }
    \Bigg],
\end{multline}
where $\phi_v$ is the {\it Hartree} term which renormalises the chemical potential, allowing for a Peirl's instability. 
$\Phi_v$ is a {\it Fock} term, describing renormalistion of the hopping parameter $v$.
The additional {\it Gorkov} terms $\Delta_v$ capture fermion-fermion pairing, relevant to theories of magnetism and superconductivity.
These terms are known as the method of {\it (superconducting) density functional theory, SC-DFT}. 
The self-consistent solution to the mean field is the saddle point 
$\langle\hat{\mathcal{H}}_{ I }\rangle = \langle\hat{\mathcal{H}}^{'}_{ I }\rangle$.
The solutions are
\begin{eqnarray}
    \rho_{m} &=& (-U_v+U_w))\langle\hat{n}^{ (m^{'}) }_{ x,y }\rangle,\\
    \Phi_v &=& -U_v\langle\hat{c}^{ (A) \dagger }_{ x,y }\hat{c}^{ (B) }_{ x,y }\rangle,\\
    \Phi_w &=& U_w\langle\hat{c}^{ (A) \dagger }_{ x,y }\hat{c}^{ (B) }_{ x-1,y }\rangle,\\
    \Delta_v &=& +U_v\langle\hat{c}^{ (A) }_{ x,y }\hat{c}^{ (B) }_{ x,y }\rangle,\\
    \Delta_w &=& -U_w\langle\hat{c}^{ (A) }_{ x,y }\hat{c}^{ (B) }_{ x-1,y }\rangle,\\
    \Delta_v^* &=& +U_v\langle \hat{c}^{ (A) \dagger}_{ x,y }\hat{c}^{ (B) \dagger }_{ x,y }\rangle,\\
    \Delta_w^* &=& -U_w\langle \hat{c}^{ (A) \dagger}_{ x,y }\hat{c}^{ (B) \dagger }_{ x-1,y }\rangle,
\end{eqnarray}
where $m'$ refers to the opposite atom to $m$. 
Note the sign of the Gorkov terms due Fermionic anticommutation.

For simplicity and because we are in particular interested in the multiorbital pairing due to $U_v$, we neglect the intercell interaction $U_w$.

\section{Results with impurities}
The impurity is located at the centre with coupling strength $V=1.21$.

\createfigure{2D-Weyl-SSH_impurity/ldos}{2D-SSH_ldos}
\createfigure{2D-Weyl-SSH_impurity/k_space}{2D-SSH_k_space}
\createfigure{2D-Weyl-SSH_impurity/majorana_fermi_arc}{majorana_fermi_arc}

\createfigure{2D-Weyl-SSH_impurity_wall/ldos}{2D-SSH_ldos}
\createfigure{2D-Weyl-SSH_impurity_wall/k_space}{2D-SSH_k_space}
\createfigure{2D-Weyl-SSH_impurity_wall/majorana_fermi_arc}{majorana_fermi_arc_impurity}

The $x$-axis is glued to periodic boundary conditions.
\createfigure{2D-Weyl-SSH_impurity_wall_PBC/majorana_fermi_arc}{majorana_fermi_arc_impurity_PBC}

Our investigation shows no interesting interplay between a single impurity and topological modes,
however, impurities can be used to generate quasiparticle-interference patterns as an experimental diagnosis of topological modes.
Moreover, a wall of impurities 

\section{Quasiparticle-interference patterns for Majorana and Bogoliubov-Fermi arcs}

\createfigure{2D-Weyl-SSH_impurity/fermi_surface}{fermi_surface}

\createfigure{2D-Weyl-SSH_impurity/qpi}{qpi}

\section{Phase diagram}

\createfigure{2D-Weyl-SSH/phase_diagram_closed}{weyl_phase_diagram_closed}
\createfigure{2D-Weyl-SSH/phase_diagram_open}{weyl_phase_diagram_open}

\section{Wall impurities}

\createfigure{2D-Weyl-SSH-impurities/phase_diagram_closed}{wall_phase_diagram_closed}
\createfigure{2D-Weyl-SSH-impurities/phase_diagram_open}{wall_phase_diagram_open}

\createfigure{2D-Weyl-SSH-impurities/V_phase_diagram_topological}{V_phase_diagram_topological}
\createfigure{2D-Weyl-SSH-impurities/V_phase_diagram_transition}{V_phase_diagram_transition}
\createfigure{2D-Weyl-SSH-impurities/V_phase_diagram_trivial}{V_phase_diagram_trivial}
